\section{Метрика LeaderRank}
\label{sec:metric}

\subsection{Что это и как работает}
LeaderRank --- метрика центральности, то есть она даёт каждой вершине одно число
<<важности>>. Это модификация PageRank, придуманная для поиска лидеров в социальных
сетях. Главное отличие от PageRank в том, что у неё нет настраиваемых параметров.

Идея простая. В граф добавляют одну искусственную вершину --- \emph{фоновый узел} $g$ (ground node).
Его соединяют двусторонними рёбрами со всеми остальными вершинами ($g\to i$ и $i\to g$ для каждой
вершины $i$). По такому расширенному графу из $n+1$ вершины запускают обычное случайное блуждание,
без телепортации и без коэффициента затухания.

Почему это работает. Фоновый узел связан со всеми в обе стороны, поэтому расширенный граф сильно
связный и апериодичный. Из теории цепей Маркова отсюда следует, что у блуждания есть единственное
стационарное распределение, а степенной метод к нему сходится. И всё это без подбора параметров.
Заодно фоновый узел играет ту же роль, что телепортация в PageRank, но встроен прямо в структуру
графа.

\subsection{Формула}
Пусть у вершины $i$ исходящая степень $\degout(i)$. В расширенном графе добавляется ребро к фоновому
узлу, поэтому исходящая степень становится $\degout(i)+1$. На каждом шаге каждая вершина $i$ отдаёт
по $s_i/(\degout(i)+1)$ каждому реальному соседу и столько же фоновому узлу, а фоновый узел отдаёт по
$s_g/n$ каждой из $n$ вершин. Начальное состояние: $s_i=1$ у всех реальных вершин, $s_g=0$. Суммарная
<<масса>> равна $n$ и сохраняется на всех шагах. После сходимости масса фонового узла поровну
возвращается вершинам: $s_i \mathrel{+}= s_g/n$.

Удобно ввести вклад вершины, который считается раз за итерацию:
\[
  \mathrm{contrib}[j] = \frac{s_{\mathrm{old}}[j]}{\degout(j)+1}.
\]
Тогда обновление состояния записывается так:
\begin{align}
  \mathrm{base} &= s_{\mathrm{old}}[g]/n, \label{eq:base}\\
  s_{\mathrm{new}}[i] &= \mathrm{base} + \sum_{j\,:\,j\to i}\mathrm{contrib}[j],
      \qquad i=1,\dots,n, \label{eq:pull}\\
  s_{\mathrm{new}}[g] &= \sum_{j=1}^{n}\mathrm{contrib}[j]. \label{eq:ground}
\end{align}
Формулы~\eqref{eq:base} и~\eqref{eq:ground} --- дешёвые операции порядка $\bigO(n)$. Дорогая часть
одна --- сумма~\eqref{eq:pull} по рёбрам для каждого приёмника. В матричном виде это
$s_{\mathrm{new}} = \mathrm{base}\cdot\one + A\T\cdot\mathrm{contrib}$, то есть умножение
транспонированной матрицы смежности на вектор. Важно, что $2n$ рёбер фонового узла физически не
хранятся: его вклад считается двумя формулами.

\subsection{Итерация --- это один проход по рёбрам}
Ключевой факт: умножение разреженной матрицы на вектор --- это один проход по списку рёбер. Чтобы
посчитать $A\T\cdot\mathrm{contrib}$, достаточно для каждого ребра $j\to i$ прибавить
$\mathrm{contrib}[j]$ к ячейке $i$. Матрицу в памяти держать не нужно, нужен только поток рёбер.
Именно на этом держится всё решение: рёбра можно читать с диска по очереди, а в памяти хранить лишь
векторы размера $\bigO(n)$.

Останавливаемся, когда состояние почти перестаёт меняться. Критерий --- норма $L_1$ разности
$\sum_i |s_{\mathrm{new}}[i] - s_{\mathrm{old}}[i]| < \varepsilon$. Про конкретный порог и число
итераций --- в разделе~\ref{sec:complexity}. Масса сохраняется, поэтому нормировать вектор на каждом
шаге не нужно. Редкая перенормировка к сумме $n$ применяется только чтобы не накапливалась ошибка
округления.

\subsection{Где это полезно}
LeaderRank ищет влиятельные узлы. Исходно его применяли к социальной сети Delicious для поиска
лидеров мнений. По сути он годится везде, где применим PageRank: ранжирование
страниц, важность пользователей, графы цитирования, рекомендации. Его берут, когда не хочется
подбирать коэффициент затухания или когда в графе много висячих вершин.

\subsection{Почему выбрана именно она}
\label{sec:why-leaderrank}
Условие разрешает любую метрику, применяемую ко всему графу, и рекомендует PageRank, LeaderRank,
Katz, Jaccard, Dice. LeaderRank выбран по нескольким причинам.
\begin{enumerate}[itemsep=3pt]
  \item \textbf{Подходит формат.} LeaderRank даёт одно число на вершину, что прямо ложится в вывод
        \texttt{vertex,rank}. Пример выходного файла в условии так и называется ---
        \texttt{leader\_rank.csv}.
  \item \textbf{Естественно считается с диска.} Каждая итерация --- один последовательный проход по
        рёбрам. Значит рёбра (это и есть тот <<гигабайт>>) читаются потоком с диска, а в памяти
        остаются только векторы $\bigO(n)$. Память ограничена числом вершин, а не числом рёбер.
  \item \textbf{Не боится гипер-узлов.} Списки смежности не материализуются в памяти. Гипер-узел
        степени $50\,000$ --- это просто $50\,000$ записей в потоке рёбер на диске. Вклад вершины
        делится на её степень, никакого взрыва нет.
  \item \textbf{Легко параллелится.} Проход по рёбрам просто делится между потоками.
  \item \textbf{Детерминирован.} Случайности в алгоритме нет. При фиксированном порядке сложения
        результат повторяется.
  \item \textbf{Выигрывает у соседей по семейству.} В отличие от PageRank, не нужен параметр $d$ и
        не нужна отдельная обработка висячих вершин --- её берёт на себя фоновый узел. В отличие от
        Katz, не нужен спектральный радиус. Меры Jaccard и Dice вообще не подходят: требуют пересечения списков смежности (случайный доступ --- худший случай
        для диска) и катастрофически взрываются на гипер-узлах (узел степени $D$ участвует в
        $\bigO(D^2)$ парах). При этом схема вычислений у LeaderRank та же, что у PageRank, поэтому при
        желании PageRank поддерживается почти тем же кодом.
\end{enumerate}
